% !TEX root = ../matan.tex

\section{Ряды Тейлора, определение $e^z$, $\cos z$, $\sin z$.}

\begin{Definition}{Ряд Тейлора}{}
	Пусть функция $S$ бесконечно дифференцируема в точке $x_0$. Её \textbf{рядом Тейлора} (с центром $x_0$) называется степенной ряд
	\[
	\sum_{k=0}^{\infty} \frac{S^{(k)}(x_0)}{k!}\,(x - x_0)^k.
	\]
	Если $S$ аналитична в $x_0$, то ряд Тейлора сходится к $S$ в окрестности этой точки (см. вопросы~32--33); в общем случае класса $C^\infty$ он может сходиться не к $S$.
\end{Definition}

\begin{Definition}{Показательная и тригонометрические функции}{}
	Для всех $z \in \C$ определяются степенными рядами
	\[
	e^z := \sum_{k=0}^{\infty} \frac{z^k}{k!}, \qquad
	\cos z := \sum_{k=0}^{\infty} (-1)^k \frac{z^{2k}}{(2k)!}, \qquad
	\sin z := \sum_{k=0}^{\infty} (-1)^k \frac{z^{2k+1}}{(2k+1)!}.
	\]
\end{Definition}

\begin{Statement}{Сходимость на всей плоскости}{}
	Все три ряда сходятся абсолютно при любом $z \in \C$ (радиус сходимости $R = +\infty$).
\end{Statement}

\begin{proof}
	Для ряда $e^z$ применим признак Даламбера к ряду модулей: при $z \neq 0$
	\[
	\frac{|z|^{k+1}/(k+1)!}{|z|^{k}/k!} = \frac{|z|}{k + 1} \xrightarrow{k \to \infty} 0 < 1,
	\]
	поэтому ряд сходится абсолютно при всех $z$, то есть $R = +\infty$. Ряды для $\cos z$ и $\sin z$ получаются из ряда $e^z$ выбором части членов с чередованием знаков, поэтому мажорируются рядом $\sum |z|^k/k!$ и также сходятся абсолютно при всех $z$.
\end{proof}

\begin{Statement}{Формула Эйлера}{}
	Для всех $z \in \C$
	\[
	e^{iz} = \cos z + i \sin z.
	\]
	В частности, при вещественном $\varphi$: $e^{i\varphi} = \cos\varphi + i\sin\varphi$, откуда
	\[
	\cos z = \frac{e^{iz} + e^{-iz}}{2}, \qquad \sin z = \frac{e^{iz} - e^{-iz}}{2i}.
	\]
\end{Statement}

\begin{proof}
	Подставим $iz$ в ряд для $e^{w}$ и сгруппируем члены по чётности (перестановка членов законна ввиду абсолютной сходимости). Используя $i^{2k} = (-1)^k$ и $i^{2k+1} = (-1)^k i$,
	\[
	e^{iz} = \sum_{n=0}^{\infty} \frac{(iz)^n}{n!}
	= \sum_{k=0}^{\infty} \frac{(iz)^{2k}}{(2k)!} + \sum_{k=0}^{\infty} \frac{(iz)^{2k+1}}{(2k+1)!}
	= \sum_{k=0}^{\infty} (-1)^k \frac{z^{2k}}{(2k)!} + i\sum_{k=0}^{\infty} (-1)^k \frac{z^{2k+1}}{(2k+1)!} = \cos z + i\sin z.
	\]
	Заменяя $z$ на $-z$, получаем $e^{-iz} = \cos z - i\sin z$ (так как $\cos$ чётна, $\sin$ нечётна); складывая и вычитая два равенства, выражаем $\cos z$ и $\sin z$.
\end{proof}

\begin{Remark}{}{}
	Поскольку $R = +\infty$, функции $e^z, \cos z, \sin z$ аналитичны на всей плоскости, и их ряды суть ряды Тейлора с центром в нуле. Для вещественного $x$ эти определения совпадают с обычными $e^x$, $\cos x$, $\sin x$, а основное свойство $e^{z+w} = e^z e^w$ проверяется перемножением рядов.
\end{Remark}
